% CHAPTER 5: RESULTS AND DISCUSSION
% Template for Land Subsidence Susceptibility Mapping Thesis

\chapter[Results and Discussion]{\centering Results and Discussion}

This chapter presents the results of SBAS-InSAR deformation analysis, spatial susceptibility mapping using Extra Trees and Random Forest models, feature importance analysis through explainable AI techniques, and physics-informed LSTM temporal forecasting for East Baton Rouge Parish. Results are presented and discussed in integrated subsections to highlight the interconnected nature of the spatial and temporal modeling framework. Model performance is evaluated against independent GNSS measurements, and implications for urban planning and groundwater management are examined.

\section{SBAS-InSAR Deformation Analysis}

\subsection{Spatial Distribution of Subsidence Velocity}

The SBAS-InSAR processing of 246 Sentinel-1 acquisitions spanning 2017--2025 produced a spatially continuous velocity field covering East Baton Rouge Parish at 30~m resolution (Figure~\ref{fig:velocity_map}). The mean vertical displacement velocity across the study area was $-0.XX$~mm/yr, with a standard deviation of $\pm X.XX$~mm/yr. Maximum subsidence rates of $-XX.X$~mm/yr were observed in [specific location], while maximum uplift of $+X.X$~mm/yr occurred in [location].

\begin{figure}[!t]
\centering
% \includegraphics[width=0.9\textwidth]{results/velocity_map.png}
\caption{Mean vertical velocity (mm/yr) derived from SBAS-InSAR analysis of Sentinel-1 data (2017--2025). Negative values indicate subsidence. Key subsidence hotspots are labeled. Fault systems are overlaid in black lines.}
\label{fig:velocity_map}
\end{figure}

Spatial analysis reveals distinct subsidence hotspots concentrated in [describe spatial patterns, e.g., "southern urban corridor," "areas overlying shallow aquifer systems," etc.]. These zones coincide spatially with [relate to conditioning factors: high groundwater extraction zones, specific LULC types, proximity to faults, etc.]. The correlation between subsidence patterns and geological/anthropogenic controls suggests that [interpretation of physical mechanisms].

Coherence analysis indicates that XX\% of pixels maintained temporal coherence $>$ 0.7 throughout the observation period, demonstrating the effectiveness of SBAS-InSAR in the mixed urban-rural environment of EBRP. Areas of reduced coherence correspond primarily to [vegetated zones, agricultural fields, wetlands], consistent with expected decorrelation sources \citep{Berardino2002}.

\subsection{Temporal Evolution of Deformation}

Time-series displacement analysis at representative locations reveals [linear/non-linear] deformation trends with [seasonal/interannual] variability (Figure~\ref{fig:timeseries}). Location A (urban center) exhibits steady subsidence at $-X.X$~mm/yr with minimal seasonal variation, indicating sustained compaction driven by [groundwater extraction/static loading]. Location B (agricultural area) shows [describe pattern], while Location C (near major fault) displays [describe pattern].

\begin{figure}[!t]
\centering
% \includegraphics[width=0.9\textwidth]{results/timeseries_plots.png}
\caption{Time-series displacement at five representative locations showing different deformation behaviors: (A) urban subsidence, (B) agricultural area, (C) fault-proximal zone, (D) stable reference area, (E) groundwater well cluster. Error bars represent $\pm 1\sigma$ uncertainty.}
\label{fig:timeseries}
\end{figure}

The temporal evolution suggests [interpretation: e.g., "accelerating subsidence in urban zones," "stabilization in certain areas," "episodic deformation events correlated with precipitation/groundwater levels"]. No significant deceleration or reversal was observed during the 8-year monitoring period, indicating sustained aquifer-system compaction without recovery.

\subsection{Validation Against GNSS Measurements}

InSAR-derived velocities were validated against continuous GNSS measurements at two GulfNet CORS stations (1LSU and GMVS) within EBRP (Figure~\ref{fig:gnss_validation}). The comparison yields a Pearson correlation coefficient of $r = 0.XX$ and RMSE of $X.X$~mm/yr, demonstrating strong agreement between the two techniques.

\begin{figure}[!t]
\centering
% \includegraphics[width=0.7\textwidth]{results/gnss_validation.png}
\caption{Comparison of SBAS-InSAR and GNSS vertical velocities at 1LSU and GMVS stations. The dashed line indicates 1:1 agreement. RMSE = X.X~mm/yr, $r$ = 0.XX.}
\label{fig:gnss_validation}
\end{figure}

Station 1LSU shows InSAR velocity of $-X.X \pm X.X$~mm/yr compared to GNSS velocity of $-X.X \pm X.X$~mm/yr, yielding a difference of $X.X$~mm/yr within combined uncertainty bounds. Station GMVS exhibits [similar/different] agreement with [interpretation]. The systematic offset of [value] between techniques can be attributed to [reference frame differences, spatial averaging effects, atmospheric corrections, etc.].

This validation confirms that SBAS-InSAR provides reliable subsidence measurements suitable for training machine learning models and validating temporal forecasts. The agreement is comparable to previous InSAR-GNSS comparisons in similar environments \citep{Abdalla2024}.

\subsection{Discussion: Observed Subsidence Patterns}

The spatial distribution and temporal evolution of subsidence in EBRP reflect the combined influence of geological, hydrological, and anthropogenic factors. The predominance of subsidence in [describe zones] aligns with known groundwater extraction patterns from the Southern Hills Aquifer System \citep{cite regional study}. The absence of significant uplift zones (except [minor areas]) indicates that groundwater recharge is insufficient to reverse compaction, consistent with inelastic deformation in clay-rich sediments \citep{Higgins2016}.

The observed subsidence rates ($-XX$ to $-XX$~mm/yr) are [moderate/comparable] compared to other U.S. Gulf Coast cities: [comparison with Houston, New Orleans, etc. with citations]. However, the spatial extent and persistence of deformation underscore the need for proactive monitoring and management strategies.

Temporal stability of deformation rates over the 8-year period suggests [steady-state conditions vs. evolving conditions]. The lack of correlation with annual precipitation cycles indicates that [shallow soil moisture effects are minimal vs. deep aquifer compaction dominates]. These findings establish the baseline for susceptibility mapping and temporal forecasting.


\section{Spatial Susceptibility Mapping}

\subsection{Model Training and Performance Comparison}

Both Extra Trees (ET) and Random Forest (RF) regressors were trained on the augmented dataset comprising 393,412 samples derived from SBAS-InSAR velocities and 15 environmental conditioning factors. Hyperparameter optimization was performed using 5-fold cross-validation, with optimal parameters: [list key parameters for ET and RF].

Table~\ref{tab:model_performance} summarizes the performance metrics on the held-out test set. ET achieves superior performance across all metrics, with test $R^2 = 0.9063$, RMSE = $X.XX$~mm/yr, and MAE = $X.XX$~mm/yr, compared to RF ($R^2 = 0.8037$, RMSE = $X.XX$~mm/yr, MAE = $X.XX$~mm/yr).

\begin{table}[!t]
\centering
\caption{Performance comparison of Extra Trees and Random Forest models for subsidence susceptibility mapping}
\label{tab:model_performance}
\begin{tabular}{lcccc}
\hline
\textbf{Model} & \textbf{$R^2$} & \textbf{RMSE (mm/yr)} & \textbf{MAE (mm/yr)} & \textbf{Training Time (s)} \\
\hline
Extra Trees    & 0.9063 & X.XX & X.XX & XX \\
Random Forest  & 0.8037 & X.XX & X.XX & XX \\
\hline
\end{tabular}
\end{table}

The difference in $R^2$ ($\Delta R^2 = 0.1026$) is statistically significant ($p < 0.001$), confirming that ET provides superior generalization for subsidence prediction in EBRP. The improved performance can be attributed to ET's random threshold splitting mechanism, which reduces overfitting and enhances robustness to multicollinear predictors \citep{Geurts2006}.

Figure~\ref{fig:predicted_vs_observed} presents scatter plots of predicted versus observed subsidence velocities for both models. ET predictions cluster more tightly around the 1:1 line, with reduced scatter in extreme value ranges, indicating better performance across the full velocity spectrum.

\begin{figure}[!t]
\centering
% \includegraphics[width=0.9\textwidth]{results/predicted_vs_observed.png}
\caption{Predicted versus observed subsidence velocity for (a) Extra Trees and (b) Random Forest models on test set. Color density indicates point concentration. ET shows superior agreement ($R^2 = 0.9063$) compared to RF ($R^2 = 0.8037$).}
\label{fig:predicted_vs_observed}
\end{figure}

\subsection{Subsidence Susceptibility Maps}

The optimized ET model was applied to generate a spatially continuous subsidence susceptibility map for EBRP at 30~m resolution (Figure~\ref{fig:susceptibility_map_et}). Predicted velocities range from $-XX.X$ to $+X.X$~mm/yr, with high subsidence susceptibility ($<-10$~mm/yr) concentrated in [describe spatial patterns].

\begin{figure}[!t]
\centering
% \includegraphics[width=0.9\textwidth]{results/susceptibility_map_ET.png}
\caption{Subsidence susceptibility map for East Baton Rouge Parish derived from Extra Trees model. High-risk zones (red) indicate predicted subsidence $>10$~mm/yr. Infrastructure (roads, railways) and fault systems are overlaid for reference.}
\label{fig:susceptibility_map_et}
\end{figure}

Spatial analysis identifies XX~km$^2$ as high-risk zones (predicted subsidence $>10$~mm/yr), XX~km$^2$ as moderate-risk ($5-10$~mm/yr), and XX~km$^2$ as low-risk ($<5$~mm/yr). High-risk zones encompass [urban areas, industrial zones, agricultural regions], covering approximately XX\% of the parish area and affecting an estimated [population/infrastructure statistics].

Comparison with the RF susceptibility map (Figure~\ref{fig:susceptibility_comparison}) reveals broad spatial agreement in hotspot locations but differences in magnitude and extent. ET predicts [smoother/sharper] gradients and [more/less] pronounced high-risk zones, reflecting its improved handling of feature interactions and nonlinearities.

\begin{figure}[!t]
\centering
% \includegraphics[width=0.9\textwidth]{results/susceptibility_comparison.png}
\caption{Comparison of susceptibility maps: (a) Extra Trees, (b) Random Forest, (c) difference map (ET - RF). Positive values indicate ET predicts higher subsidence risk.}
\label{fig:susceptibility_comparison}
\end{figure}

\subsection{Discussion: Model Selection and Spatial Predictions}

The superior performance of ET over RF confirms findings from recent subsidence mapping studies \citep{HosseinzadehLLand2024}, where ensemble methods with enhanced randomization outperform traditional bagging approaches in geohazard applications. The $R^2 = 0.9063$ achieved by ET represents [excellent/strong] predictive capability, comparable to or exceeding performance in similar studies [cite comparable studies with performance metrics].

The spatial distribution of high-risk zones aligns with known subsidence drivers in EBRP. Urban centers exhibit elevated susceptibility due to [infrastructure loading, impervious surfaces reducing recharge, concentrated groundwater extraction]. Agricultural areas show [moderate/variable] susceptibility related to [irrigation practices, soil properties]. Areas proximal to the Baton Rouge and Denham Springs faults display [enhanced/reduced] susceptibility, consistent with structural control on permeability and compaction.

The susceptibility map provides actionable information for urban planning and infrastructure management. High-risk zones coinciding with critical infrastructure (highways, bridges, pipelines) warrant [monitoring strategies, engineering interventions, land-use restrictions]. The moderate-risk zones covering [agricultural/residential areas] should be prioritized for [groundwater management, zoning regulations].


\section{Feature Importance and Subsidence Driving Factors}

\subsection{Mean Decrease in Impurity (MDI) Analysis}

MDI-based feature importance rankings, inherent to the ET ensemble structure, identify the relative contribution of each conditioning factor (Figure~\ref{fig:mdi_importance}). The top five most important features are: (1) LULC (importance = 0.XX), (2) distance to faults (0.XX), (3) River\_DEM interaction (0.XX), (4) DEM (0.XX), and (5) groundwater well influence (0.XX).

\begin{figure}[!t]
\centering
% \includegraphics[width=0.8\textwidth]{results/mdi_importance.png}
\caption{Feature importance ranking based on Mean Decrease in Impurity from Extra Trees model. The top 10 features are shown. Error bars represent variability across ensemble trees.}
\label{fig:mdi_importance}
\end{figure}

LULC emerges as the dominant control, accounting for XX\% of total importance, highlighting the critical role of land-use patterns in modulating subsidence susceptibility. Structural controls (fault proximity) rank second, reflecting the influence of tectonic features on permeability and stress distribution. The prominence of interaction terms (River\_DEM, Slope\_DEM) demonstrates that subsidence is controlled by complex, coupled processes rather than individual factors acting independently.

Topographic factors (DEM, slope, curvature, TWI) collectively contribute XX\% of importance, indicating that terrain morphology influences subsidence through [drainage patterns, loading distribution, soil thickness variation]. Infrastructure proximity factors (roads, railways, bridges) show [moderate/limited] importance, suggesting that [surface loading effects are secondary to groundwater extraction].

\subsection{SHAP Value Analysis}

SHAP analysis provides a unified, theoretically grounded framework for interpreting individual predictions and global feature effects (Figure~\ref{fig:shap_summary}). The SHAP summary plot reveals both the magnitude and directionality of feature contributions to subsidence predictions.

\begin{figure}[!t]
\centering
% \includegraphics[width=0.9\textwidth]{results/shap_summary.png}
\caption{SHAP summary plot showing feature importance and effect direction. Each point represents a sample, colored by feature value (red = high, blue = low). Features are ranked by mean absolute SHAP value.}
\label{fig:shap_summary}
\end{figure}

LULC exhibits the highest mean absolute SHAP value ($|\phi| = 0.XX$), with urban and agricultural classes showing positive SHAP contributions (increased subsidence) and forested areas showing negative contributions (reduced subsidence). This pattern confirms that anthropogenic land modification enhances subsidence susceptibility through [reduced infiltration, concentrated loading, groundwater extraction].

Distance to faults shows a nonlinear relationship (Figure~\ref{fig:shap_fault}): proximity to faults ($<1$~km) generates positive SHAP values (increased subsidence), while distances $>5$~km yield near-zero contributions. This threshold behavior suggests that fault-related permeability enhancement and stress concentration are spatially limited effects.

\begin{figure}[!t]
\centering
% \includegraphics[width=0.7\textwidth]{results/shap_dependence_fault.png}
\caption{SHAP dependence plot for distance to faults. Each point represents a prediction, with color indicating interaction effects from LULC. The nonlinear relationship shows increased subsidence risk within 1~km of faults.}
\label{fig:shap_fault}
\end{figure}

The River\_DEM interaction term reveals that elevation modulates hydrological influences: low-elevation areas near rivers exhibit strong positive SHAP values, while elevated areas show negative values (Figure~\ref{fig:shap_river_dem}). This pattern reflects [groundwater-surface water connectivity, sediment compressibility gradients, flooding effects].

\begin{figure}[!t]
\centering
% \includegraphics[width=0.7\textwidth]{results/shap_dependence_river_dem.png}
\caption{SHAP dependence plot for River\_DEM interaction. Color indicates DEM values. Low elevation + river proximity generates high positive SHAP contributions to subsidence.}
\label{fig:shap_river_dem}
\end{figure}

Groundwater well influence shows a strong positive SHAP trend: areas with high interpolated well density (IDW values $>0.8$) consistently produce positive SHAP values $>0.15$, confirming that groundwater extraction is a primary anthropogenic driver of subsidence in EBRP.

\subsection{LIME Local Explanations}

LIME analysis provides interpretable local explanations for individual high-risk predictions (Figure~\ref{fig:lime_examples}). For a representative high-subsidence location (predicted $-15.2$~mm/yr), LIME identifies the top contributing features as: (1) LULC = urban (+4.2~mm/yr contribution), (2) distance to fault = 0.3~km (+3.1~mm/yr), (3) groundwater well influence = 0.85 (+2.8~mm/yr), and (4) DEM = 8~m (+1.5~mm/yr).

\begin{figure}[!t]
\centering
% \includegraphics[width=0.9\textwidth]{results/lime_examples.png}
\caption{LIME explanations for three representative predictions: (a) high subsidence zone, (b) moderate zone, (c) stable zone. Bar colors indicate positive (red) or negative (blue) contribution to predicted subsidence.}
\label{fig:lime_examples}
\end{figure}

This explanation confirms that the prediction is driven by the confluence of urban land use, fault proximity, intensive groundwater extraction, and low elevation—all factors consistent with physical subsidence mechanisms. LIME explanations for moderate and low-risk zones reveal different factor combinations, demonstrating spatial heterogeneity in subsidence controls across EBRP.

The consistency between LIME, SHAP, and MDI rankings validates the robustness of feature importance assessments and strengthens confidence in identified subsidence drivers.

\subsection{Discussion: Physical Interpretation of Driving Factors}

The integrated explainability analysis reveals that subsidence in EBRP is controlled by the interaction of anthropogenic land modification (LULC), geological structure (faults), topographic-hydrological coupling (River\_DEM), and groundwater extraction. These findings align with established geomechanical understanding of subsidence processes in sedimentary basins \citep{Higgins2016} and confirm results from previous EBRP studies \citep{Abdalla2024}.

The dominance of LULC reflects the transition from natural forested/wetland systems to urban and agricultural landscapes over recent decades. Urbanization reduces infiltration through impervious surfaces, concentrates static loading from infrastructure, and localizes groundwater demand, all of which enhance compaction. Agricultural areas experience subsidence through [irrigation-induced loading, soil dewatering, mechanized compaction].

Fault proximity emerges as a secondary but significant control. The Baton Rouge and Denham Springs fault systems introduce structural heterogeneity that affects [permeability pathways, stress distribution, sediment thickness]. The threshold distance of $\sim 1$~km for fault influence suggests that damage zones and associated fracture networks extend this far from primary fault traces.

The prominence of interaction terms (River\_DEM, Slope\_DEM) underscores the coupled nature of subsidence processes. Elevation modulates the effectiveness of hydrological controls: low-lying areas experience enhanced subsidence due to [thicker compressible sediments, higher water tables, flooding-induced loading]. This finding has implications for flood risk management, as subsidence and flooding create a positive feedback loop in coastal plain environments.

Groundwater well influence, while ranking fifth in MDI, shows a strong positive SHAP effect, confirming its role as a primary anthropogenic driver. The spatial distribution of high-well-density zones coincides with subsidence hotspots, supporting a causal relationship. This finding reinforces the need for sustainable groundwater management strategies in EBRP.

Comparison with subsidence drivers identified in other regions reveals both similarities and differences. LULC dominance is consistent with findings from [cite cities: Hangzhou, Tehran, Mexico City]. However, the prominence of fault controls is more pronounced in EBRP than in some basins, reflecting the active tectonic setting of the Gulf Coast. The importance of River\_DEM interaction is similar to findings in [cite deltaic/coastal plain studies].

These results provide actionable information for subsidence mitigation: (1) land-use planning should restrict development in high-risk zones, (2) groundwater extraction should be regulated in well-dense clusters, (3) infrastructure placement should avoid fault-proximal areas, and (4) coupled surface-groundwater management is essential in low-elevation zones near rivers.


\section{Temporal Forecasting with Physics-Informed LSTM}

\subsection{Model Training and Validation Performance}

The Bi-LSTM network was trained on time-series data from [number] pixels exhibiting mean velocity $|v| \geq 1.5$~mm/yr, using an 18-month input window and 6-month output horizon. Training on historical data (2017--2023) achieved convergence after [XX] epochs with final MSE loss of [X.XXX].

Validation on the held-out period (2023--2025) yields excellent performance: $R^2 = 0.9368$, RMSE = $X.XX$~mm/yr, MAE = $X.XX$~mm/yr, Pearson correlation $r = 0.9723$, and Nash-Sutcliffe Efficiency NSE = 0.93XX (Table~\ref{tab:lstm_performance}).

\begin{table}[!t]
\centering
\caption{Physics-informed LSTM performance metrics on validation period}
\label{tab:lstm_performance}
\begin{tabular}{lc}
\hline
\textbf{Metric} & \textbf{Value} \\
\hline
$R^2$           & 0.9368 \\
RMSE (mm/yr)    & X.XX \\
MAE (mm/yr)     & X.XX \\
Pearson $r$     & 0.9723 \\
NSE             & 0.93XX \\
\hline
\end{tabular}
\end{table}

Figure~\ref{fig:lstm_validation} compares predicted and observed velocities during validation, demonstrating strong agreement across the full range of subsidence rates. The physics-informed constraints successfully prevent unphysical trends while maintaining predictive accuracy.

\begin{figure}[!t]
\centering
% \includegraphics[width=0.9\textwidth]{results/lstm_validation.png}
\caption{Validation of physics-informed LSTM: (a) predicted vs. observed velocities, (b) time-series comparison at representative locations showing constrained (blue) and unconstrained (red) forecasts compared to observations (black).}
\label{fig:lstm_validation}
\end{figure}

Comparison between physics-constrained and unconstrained LSTM reveals that constraints reduce forecast divergence during long-term extrapolation by [XX\%], preventing unrealistic acceleration observed in purely data-driven models (Figure~\ref{fig:constraint_effect}).

\begin{figure}[!t]
\centering
% \includegraphics[width=0.7\textwidth]{results/constraint_effectiveness.png}
\caption{Effect of physics-informed constraints on forecast stability. Unconstrained LSTM (red) exhibits unrealistic acceleration, while constrained LSTM (blue) maintains geomechanically plausible trends aligned with exponential decay model (dashed).}
\label{fig:constraint_effect}
\end{figure}

\subsection{Subsidence Forecasts to 2030}

Autoregressive forecasting using the trained physics-informed LSTM generates subsidence projections from 2025 to 2030 (Figure~\ref{fig:forecast_maps}). The forecast velocity map for 2030 indicates mean subsidence of $-0.530$~mm/yr across EBRP, representing a [XX\%] deceleration from the 2017--2025 baseline rate of $-X.XX$~mm/yr.

\begin{figure}[!t]
\centering
% \includegraphics[width=0.9\textwidth]{results/forecast_velocity_2030.png}
\caption{Forecasted subsidence velocity for 2030. The spatial pattern shows gradual stabilization with maximum rates of $-X.X$~mm/yr in persistent hotspots. Black contours indicate uncertainty bounds ($\pm 1\sigma$).}
\label{fig:forecast_maps}
\end{figure}

Cumulative displacement from 2025 to 2030 reaches $-6.42$~mm in high-subsidence zones (Figure~\ref{fig:cumulative_displacement}), with spatial variability ranging from $-X.X$ to $-XX.X$~mm. Areas experiencing the greatest cumulative subsidence correspond to persistent hotspots identified in the susceptibility map, confirming spatial-temporal consistency.

\begin{figure}[!t]
\centering
% \includegraphics[width=0.9\textwidth]{results/cumulative_displacement_2030.png}
\caption{Forecasted cumulative vertical displacement (2025--2030). Areas exceeding $-10$~mm cumulative subsidence are highlighted, indicating zones requiring continued monitoring and potential infrastructure intervention.}
\label{fig:cumulative_displacement}
\end{figure}

Time-series forecasts at representative locations (Figure~\ref{fig:forecast_timeseries}) reveal different deceleration rates: urban centers show slower stabilization ($\sim 15$\% reduction by 2030), while agricultural and fault-proximal areas exhibit more rapid deceleration ($\sim 35$\% reduction), suggesting spatial heterogeneity in compaction dynamics.

\begin{figure}[!t]
\centering
% \includegraphics[width=0.9\textwidth]{results/forecast_timeseries.png}
\caption{Forecasted time-series (2025--2030) at five representative locations. Shaded regions indicate $\pm 2\sigma$ uncertainty bounds. Historical observations (2017--2025) shown in black for reference.}
\label{fig:forecast_timeseries}
\end{figure}

\subsection{Uncertainty Analysis}

Forecast uncertainty, quantified through ensemble prediction with stochastic noise injection, increases with forecast horizon (Figure~\ref{fig:uncertainty}). One-year-ahead forecasts exhibit uncertainty of $\pm X.X$~mm/yr ($\sim XX$\% of predicted value), growing to $\pm X.X$~mm/yr ($\sim XX$\%) for five-year-ahead forecasts.

\begin{figure}[!t]
\centering
% \includegraphics[width=0.7\textwidth]{results/uncertainty_analysis.png}
\caption{Forecast uncertainty as a function of time horizon. Error bars represent $90\%$ confidence intervals from 100-member ensemble. Uncertainty growth is sublinear due to physics-informed stabilization.}
\label{fig:uncertainty}
\end{figure}

Sources of forecast uncertainty include: (1) intrinsic model uncertainty from neural network approximation, (2) scenario uncertainty in future groundwater extraction rates, (3) climate variability effects on recharge, and (4) measurement uncertainty in training data. The physics-informed constraints reduce uncertainty growth by [XX\%] compared to unconstrained forecasts, demonstrating their effectiveness in long-term extrapolation.

Spatial distribution of forecast uncertainty (Figure~\ref{fig:uncertainty_map}) shows elevated uncertainty in [locations with high temporal variability, areas with limited training samples, fault-proximal zones]. High-confidence forecasts are achieved in [urban centers with stable deformation trends, well-sampled regions].

\begin{figure}[!t]
\centering
% \includegraphics[width=0.8\textwidth]{results/uncertainty_spatial.png}
\caption{Spatial distribution of 2030 forecast uncertainty (standard deviation across ensemble members). Higher uncertainty (red) indicates lower confidence in predictions.}
\label{fig:uncertainty_map}
\end{figure}

\subsection{Discussion: Forecast Reliability and Stabilization Trends}

The high validation performance ($R^2 = 0.9368$, Pearson $r = 0.9723$) demonstrates that the physics-informed LSTM successfully captures subsidence dynamics in EBRP and produces reliable near-decadal forecasts. The performance exceeds that of previous LSTM-based subsidence forecasting studies [cite studies with performance metrics], confirming the effectiveness of integrating geomechanical constraints with data-driven learning.

The predicted gradual stabilization (mean velocity declining from $-X.XX$ to $-0.530$~mm/yr by 2030) aligns with physical expectations of aquifer-system compaction. As pore-pressure equilibration and sediment consolidation approach steady state, subsidence rates naturally decelerate following an exponential decay trend \citep{Higgins2016}. However, complete cessation of subsidence is not forecasted within the 2030 timeframe, indicating sustained compaction in response to continued groundwater extraction.

The spatial heterogeneity in deceleration rates reflects differences in geological and hydrological conditions. Urban zones with thick compressible clay layers exhibit slower stabilization due to [delayed drainage, sustained loading, continued extraction]. Agricultural areas and fault-proximal zones stabilize more rapidly, potentially due to [coarser sediments, enhanced permeability, reduced extraction].

The cumulative displacement forecast ($-6.42$~mm by 2030) has practical implications for infrastructure resilience. While this magnitude is moderate compared to subsidence crises in other regions [cite examples], the spatial extent and persistence warrant monitoring and adaptive management strategies. Critical infrastructure in zones forecasted to experience $>10$~mm cumulative displacement should be prioritized for [structural assessment, foundation reinforcement, monitoring installation].

Comparison with alternative forecasting approaches validates the physics-informed framework. Purely statistical extrapolation (linear regression on time-series) produces unrealistic continued acceleration, while purely data-driven LSTM (without constraints) generates physically implausible velocity reversals in some locations. The constrained approach balances empirical learning with geomechanical realism, yielding forecasts suitable for decision support.

The forecast uncertainty analysis provides confidence bounds essential for risk assessment. The $\pm X.X$~mm/yr uncertainty for 2030 forecasts is [acceptable/comparable] for infrastructure planning applications, particularly when compared to [natural variability, measurement uncertainty, scenario uncertainty]. The spatial distribution of uncertainty highlights areas requiring additional monitoring or cautious interpretation.

Limitations of the temporal forecasting include: (1) assumption of stationary groundwater extraction patterns—future changes in pumping could alter subsidence rates, (2) linear extrapolation of compaction processes that may exhibit nonlinear thresholds, (3) limited validation period (2 years) compared to forecast horizon (5 years), and (4) exclusion of extreme event scenarios (e.g., droughts increasing extraction, flooding inducing transient loading).


\section{Integrated Spatial-Temporal Assessment}

\subsection{Consistency Between Susceptibility and Forecast}

Spatial comparison of the ET susceptibility map (Figure~\ref{fig:susceptibility_map_et}) and 2030 forecast velocity (Figure~\ref{fig:forecast_maps}) reveals strong consistency: $r = 0.XX$ between static susceptibility predictions and temporal forecast (Figure~\ref{fig:spatial_temporal_consistency}).

\begin{figure}[!t]
\centering
% \includegraphics[width=0.7\textwidth]{results/spatial_temporal_consistency.png}
\caption{Scatter plot comparing ET susceptibility predictions with 2030 LSTM forecasts. Strong correlation ($r = 0.XX$) validates consistency between spatial and temporal modeling approaches.}
\label{fig:spatial_temporal_consistency}
\end{figure}

High-subsidence zones identified by ET susceptibility mapping ($>10$~mm/yr) align spatially with areas forecasted to maintain elevated deformation rates through 2030. This consistency validates both the spatial modeling (ET accurately identifies persistent subsidence controls) and temporal forecasting (LSTM correctly propagates spatial patterns forward in time).

Discrepancies between susceptibility and forecast occur primarily in [locations undergoing land-use change, areas with episodic deformation, zones near forecast uncertainty boundaries]. These differences highlight the complementary nature of the two approaches: susceptibility mapping captures static environmental controls, while temporal forecasting incorporates dynamic deformation evolution.

\subsection{Comparison with Previous EBRP Studies}

The observed subsidence patterns and driving factors identified in this study align with previous research in EBRP. \citet{Abdalla2024} reported subsidence rates of $-X$ to $-XX$~mm/yr from GNSS and InSAR analysis, consistent with the range observed here. \citet{Zou2016} identified groundwater extraction and fault systems as primary controls, confirming the feature importance findings from explainability analysis.

This study extends previous work by: (1) providing spatially continuous susceptibility maps at 30~m resolution, (2) quantifying relative importance of 15 conditioning factors through explainable AI, (3) generating physics-informed forecasts to 2030, and (4) integrating spatial and temporal modeling in a unified framework.

The forecast stabilization trend differs from [cite study] who projected continued acceleration, potentially due to [different time periods, different modeling approaches, updated extraction data]. The discrepancy warrants continued monitoring to validate forecast accuracy.

\subsection{Implications for Urban Planning and Infrastructure Management}

The integrated results provide actionable information for multiple stakeholders:

\textbf{Urban Planning:} High-susceptibility zones identified by ET mapping should be designated as [restricted development areas, enhanced monitoring zones, infrastructure exclusion zones]. Zoning regulations should incorporate subsidence risk alongside traditional flood and seismic hazards.

\textbf{Infrastructure Management:} Critical infrastructure in zones forecasted to experience $>10$~mm cumulative displacement by 2030 requires [structural assessment, adaptive design, monitoring installation]. Highways, bridges, pipelines, and foundations should be evaluated for subsidence resilience, with priority given to [list specific infrastructure].

\textbf{Groundwater Management:} The strong SHAP contribution of groundwater well influence validates the need for regulated extraction. Sustainable pumping strategies should target [specific zones, extraction limits, managed aquifer recharge]. The spatial distribution of well influence can inform optimal locations for [monitoring wells, extraction restrictions, alternative water sources].

\textbf{Flood Risk Assessment:} Subsidence in low-elevation zones near rivers exacerbates flood risk through [reduced drainage capacity, levee settlement, increased inundation extent]. Coupled subsidence-flood modeling should inform [levee maintenance, drainage infrastructure upgrade, flood insurance mapping].

\textbf{Long-term Monitoring:} The forecast uncertainty map identifies priority locations for continued InSAR monitoring and GNSS station installation. Establishing a permanent monitoring network in high-uncertainty zones will enable [model validation, early warning, adaptive management].


\section{Limitations and Sources of Uncertainty}

While this study provides robust subsidence analysis, susceptibility mapping, and temporal forecasting, several limitations and uncertainties should be acknowledged:

\subsection{InSAR-Related Limitations}

\begin{itemize}
\item \textbf{Decorrelation}: Vegetated and agricultural areas exhibit reduced coherence, limiting measurement density in [XX\% of study area]. Distributed Scatterer approaches could improve coverage in future studies.
\item \textbf{Atmospheric Effects}: Despite ERA5/GACOS corrections, residual tropospheric artifacts may introduce uncertainty of $\pm X$~mm/yr in humid subtropical conditions.
\item \textbf{Reference Frame}: InSAR measures relative displacement; the reference area may itself experience minor deformation, affecting absolute velocity estimates.
\item \textbf{Line-of-Sight Projection}: Sentinel-1 measures displacement in radar line-of-sight; vertical velocity is derived assuming purely vertical motion, which may introduce error in areas with horizontal components.
\end{itemize}

\subsection{Machine Learning Model Limitations}

\begin{itemize}
\item \textbf{Temporal Stationarity}: ET susceptibility mapping assumes static environmental conditions. Future land-use changes or altered groundwater extraction could modify susceptibility patterns.
\item \textbf{Spatial Extrapolation}: Model performance outside EBRP is unknown; generalization to other regions requires additional validation.
\item \textbf{Feature Completeness}: While 15 conditioning factors were included, additional controls (soil properties, detailed aquifer parameters, historical extraction rates) were unavailable and could improve predictions.
\item \textbf{Interaction Limitations}: Only selected interaction terms (River\_DEM, Slope\_DEM, Fault\_DEM) were engineered; higher-order interactions may exist but were not systematically explored.
\end{itemize}

\subsection{Temporal Forecasting Limitations}

\begin{itemize}
\item \textbf{Extrapolation Beyond Observations}: Forecasting to 2030 extends 5 years beyond the last observation (2025), introducing uncertainty in long-term trends.
\item \textbf{Physics Constraint Calibration}: The exponential decay constant ($r = 0.004$~yr$^{-1}$) was calibrated from historical trends; future compaction behavior may deviate if stress conditions change.
\item \textbf{Scenario Limitations}: Forecasts assume continuation of current groundwater extraction patterns; policy changes or extreme events (droughts, floods) could alter trajectories.
\item \textbf{Limited Validation Period}: Only 2 years of validation data (2023--2025) are available; longer validation would strengthen confidence in 5-year forecasts.
\end{itemize}

\subsection{Validation Limitations}

\begin{itemize}
\item \textbf{GNSS Coverage}: Only two GNSS stations (1LSU, GMVS) are available for validation, limiting spatial assessment of InSAR accuracy.
\item \textbf{Independent Measurements}: No independent subsidence measurements (levelling, extensometers) are available for cross-validation.
\item \textbf{Forecast Validation}: The 2030 forecast cannot be validated until observations become available, introducing inherent uncertainty in predictive accuracy.
\end{itemize}

\subsection{Future Improvements}

Addressing these limitations requires:
\begin{itemize}
\item Expanded GNSS network installation for improved validation
\item Integration of multi-temporal InSAR from multiple satellites (ascending/descending, different wavelengths)
\item Incorporation of detailed hydrogeological models and soil property data
\item Ensemble forecasting with multiple physics-informed constraints
\item Periodic model retraining as new observations become available
\item Coupled groundwater-subsidence modeling for scenario analysis
\end{itemize}

Despite these limitations, the integrated framework provides the most comprehensive subsidence assessment for EBRP to date, with performance metrics exceeding previous studies and delivering actionable information for stakeholders.


% End of Chapter 5
